\chapter{Wrapping Your Host Code}
\label{ch:wrapping}
A surrogate changes one block of a coupled physical iteration. To predict
its effect on the converged state, GRAIL needs the residual of the complete
cycle and derivative actions for the two composed blocks.

\section{The Required Inputs}
\label{sec:five-things}
Supply a converged surrogate state $x_e$, the signed residual
$b=\mathcal T_e(x_e)-\mathcal T_t(x_e)$, and two derivative callbacks.
The callbacks compose as \code{emulator\_jvp(host\_jvp(v))} and must
represent the actual surrogate-containing cycle. GRAIL solves
$(I-K_e)p=b$. The estimated reference state is $x_e-p$.
The residual and both derivatives use the same operating point.
For a composition $P\circ R$, evaluate $DR$ at $x_e$ and $DP$ at
$R(x_e)$, rather than reusing the input state for both blocks.

\section{The State-Vector Seam}
\label{sec:seam}
Use one documented coordinate system throughout the residual, derivatives,
solver, and verification target. Both callbacks accept and return flat
vectors with the agreed shape and precision. Intermediate variables may
have different physical units; scaling must be consistent with the chain rule.
For conserved populations, independent coordinates or a conservation-constrained
tangent avoid redundant state directions. The structured two-level example
uses $q=n_2$ and reconstructs $n_1=1-q$.
If both populations are stored, the reference blocked convention is
\code{populations.T.flatten()}, with inverse
\code{flat.reshape(2, N\_tau).T}.

\section{Worked Example: The Structured Two-Level Closure}
\label{sec:twolevel-worked}
Run the supplied, complete example after installing the numerical extras:
\begin{lstlisting}[style=grailInStyle]
python examples/structured_response.py
\end{lstlisting}
It loads the supplied learned checkpoint, solves its coupled equilibrium,
and evaluates the residual there. The host converts populations to source
functions and transports radiation. The second block is the actual trained
neural adapter; automatic differentiation follows that adapter, not the
reference rate-equation solver. Both callbacks use the correct intermediate
radiation field. A separately solved reference equilibrium supplies the
verification target after prediction.

The executable example checks the true linear equation defect and compares
the predicted shift with the independent target. It applies the correction by
\begin{lstlisting}[style=grailInStyle]
q_reference_estimate = q_emulator - result.delta_x_star
\end{lstlisting}
The complete source shows how to construct each callback without substituting
a reference Jacobian for the neural one.

\section{Multilevel and Other Host Codes}
\label{sec:threelevel-worked}
\label{sec:your-host}
A multilevel atom changes the state dimension and constraints, not the
callback contract. A chemistry solver or a compiled host can expose an
analytic, automatic-differentiation, or independently checked finite-difference
JVP. Keep units, packing, conservation, and operating points explicit.
In electronic structure, the neural Fock construction can be the first-applied
block. The historical callback argument names do not determine which physical
block contains the neural model; composition order does.

\section{The Four Standard Pitfalls}
\label{sec:pitfalls}
\subsection{State-Vector Layout}
\label{pit:layout}
Check a pack/unpack round trip and verify the derivative on independent
coordinate directions. A reshaping error can preserve array sizes while
changing the physical operator.
\subsection{Autograd Trace Breakage}
\label{pit:autograd}
Converting to NumPy or detaching an intermediate can remove physical feedback
from the derivative. Compare the composed JVP with finite differences of the
actual unaccelerated cycle, including diagnostic closures and intermediate states.
\subsection{Dtype and Device}
\label{pit:dtype}
Use consistent float64-equivalent values at the GRAIL boundary. A smooth
float64 derivative should correspond to the same evaluator used for the cycle.
\subsection{The Neural Derivative}
\label{pit:jacobian}
Differentiate the actual replacement block. A small GMRES residual only checks
the operator supplied to GMRES; it cannot detect a consistently supplied,
incorrect reference derivative.

\section{Verification Before Deployment}
\label{sec:debug-recipe}
\label{sec:checklist}
Check the converged-state equation, the signed cycle residual, derivative
composition against an independent cycle evaluation, the true linear defect,
and physical observables. Use independently re-equilibrated states to test
local prediction accuracy. Small per-call mean-square error alone does not
control feedback derivatives or nonlinear remainders. Chapter~\ref{ch:output}
explains solver diagnostics; Chapter~\ref{ch:coverage} describes calibration.
